% Pursuing Stacks (cotes 134-2 to 134-8): archive cote 134-2, site page 26 (PDF page 27).
% letter to Quillen. Draft compared: ps1.tex:576-638.
% AI-transcribed by Claude Opus 5.5 (claude-opus-5-5, adaptive thinking, effort high), Message Batches API, batch ps1 (msgbatch_01Rv5hCaNTtUt4JHx1yqoHPL), 2026-11-12; reviewed against both supplied scan halves; see CHECK_DECISIONS.tsv.
\scanpage{134-2/26}
\margin{7}\mspage{\deleted{\illegible{1}} \hand{12}}

the boundary \xout{x} $S_i = \dot{D}_{i+1}$ of $D_{i+1}$ , more specifically these images are
just two\allowbreak{}"complementary" hemispheres in $S_i$ , which I will denote by
$S_i^+$ and $S_i^-$ . T\hand{h}e kernel of the pair $(\hand{\varphi}_i^+ , \hand{\varphi}_i^-)$ is just $S_{i-1} =$ $\dot{D}_i$ ,
and the common restriction of the maps $\hand{\varphi}_i^+ , \hand{\varphi}_i^-$ to $S_{i-1}$ is an isomorphism
\xout{xx}
\[ S_{i-1} \hand{\simeq} S_i^+ \hand{\cap} S_i^- . \]
This $S_{i-1}$ in turn decomposes into the two hemispheres $S_{i-1}^+$ , $S_{i-1}^-$ ,
images of $D_{i-1}$ . Replacing $D_{i+1}$ by $D_i$ , we see that the i-cell $D_i$
is decomposed into a union of 2i+1 closed cells, one being $D_i$ itself,
the others being canonically isomorphic to the cells $S_j^+$ , $S_j^-$ $(0 \leq j$ $\hand{<}$ $i)$,
images of $D_j \hand{\longrightarrow} D_n$ by the iterated morphisms
\[ \hand{\varphi}_{n,j}^+ , \hand{\varphi}_{n,j}^- : D_j \hand{\rightrightarrows} D_n . \]
This is a cellular decomposition, corresponding to a partition of $D_n$
into 2n+1 open cells $\mathring{D}_n$ , $\mathring{S}_j^+$ $=$ \xout{\illegible{1}} $\hand{\varphi}_{n,j}^+(\mathring{D}_j)$ , $\mathring{S}_j^-$ $= \hand{\varphi}_{n,j}^-(D_j)$ . For
any cell \inserted[typed]{in this decomposition}, the incident cells are exactly those of strictly smaller
dimension.

\margin{$\overset{n}{\underset{\ell}{*}}$}
When introducing the operation $\overset{n}{\underset{\hand{\ell}}{*}}$ with $\hand{\ell} = n-j$, it is seen that
this corresponds to choosing a map
\[ D_n \hand{\longrightarrow} (D_n , S_j^+) \hand{\amalg}_{D_j} (D_n , S_j^-) , \]
satisfying a certain condition \deleted{\unsure{$*$}}, expressing the formulas I wrote down
yesterday for $s_1$ and $t_1$ of $u \underset{\hand{\ell}}{*} v$ - the \xout{condition}\xout{x} formulas translate
\xout{\unsure{to S}}
into demanding that the restriction of the looked\hand{-}for map of $D_n$ to its
boundary $S_{n-1}$ should be a given map (given at any rate, for $\hand{\ell} \geq 2$, in
terms of the operation $\underset{\hand{\ell}-1}{*}$ ,\xout{\illegible{3}} which explains the point I
made that the $\underset{\hand{\ell}}{*}$ -structure is of order just above the $\underset{\hand{\ell}-1}{*}$ structure,
namely (inductively) is of order $\hand{\ell}$ ...). T\hand{h}at the extension of this
map of $S_{n-1}$ to $D_n$ does indeed exist, comes from the fact that the amal-%
gamated sum on the right hand side is contractible for obvious reasons.

This gives a clue of what we should call ``standar\over{t}{d}'' amalgamated sums
of the cells $D_i$ . T\hand{h}e first idea that comes to mind is that we should
insist that the space considered should be contractible, excluding
amalgamated sums therefore such as
\hand{\diagram{
\begin{tikzcd}[row sep=tiny]
 \& \bullet \ar[dr, no head, color=blue] \& \\
\bullet \ar[ur, no head, color=blue] \ar[rr, no head, color=blue] \& \& \bullet
\end{tikzcd}
\quad\text{or}\quad
\begin{tikzcd}
\bullet \ar[r, no head, bend left, color=blue] \ar[r, no head, bend right, color=blue] \& \bullet
\end{tikzcd}
}}
\inserted{, which are circles.}

This \inserted[typed]{formulation}\xout{\illegible{1}} however has the unconvenience of not being directly ex-%
pressed in combinatorial terms. T\hand{h}e following formulation, which has
the advantage of being of combinatorial nature, is presumably equivalent
to the former, and gives \inserted{at any rate}\xout{enough} (I expect) \xout{enough} a large enough notion
of ``standar\over{t}{d}ness'' to yield \over{\unsure{f}}{f}or the corresponding notion of $\infty$-category
enough structure for whatever one will ever need. In any case, it is
\margin{26}
